%==================================================================
% FORMULIR RESMI  FR-UAJ-07-39/R0
% PERNYATAAN PEMBIMBING TUGAS AKHIR MAHASISWA
%
% Data diambil otomatis dari a0-identitas.tex.
% Judul dipecah manual menjadi tiga baris karena mengikuti bentuk
% formulir resmi; sesuaikan pemenggalannya dengan judul Anda.
%==================================================================
\clearpage
\thispagestyle{empty}
\kopformulir{FR-UAJ-07-39/R0}

\begin{center}
    \textbf{FORMULIR}\\[0.3em]
    \textbf{PERNYATAAN PEMBIMBING TUGAS AKHIR MAHASISWA}\\[0.3em]
    \textbf{PROGRAM STUDI \MakeUppercase{\prodi}}
\end{center}
\vspace{1.5em}

{\bfseries
\isianform{NAMA}{\MakeUppercase{\penulis}}
\vspace{1em}
\isianform{NIM}{\nim}
\vspace{0.8em}
% --- Judul: diambil otomatis dari \judulid, membungkus sendiri ---
\isianjudul{JUDUL TUGAS AKHIR}{\judulid}
\vspace{0.8em}
\isianform{Tanggal Ujian T.A}{\tglujian}
}

\vspace{1.5em}
\noindent Menyatakan bahwa \textbf{Tugas Akhir dan Makalah} telah dikoreksi
sesuai dengan hasil sidang\textbf{.}

\vspace{2.5em}
\hfill \kotaformulir, \tglkoreksi \hspace{1.5cm}

\vspace{1em}
\noindent
\ifdefstring{\jumlahpembimbing}{2}{%
\begin{tabular}{@{}p{0.49\textwidth}@{\hfill}p{0.49\textwidth}@{}}
    Pembimbing I & Pembimbing II \\[3cm]
    {\small\bfseries\pembimbing} & {\small\bfseries\pembimbingdua} \\
    {\small NIP/NIK. \NIPpembimbing} & {\small NIP/NIK. \NIPpembimbingdua} \\
\end{tabular}
}{%
\begin{tabular}{@{}p{0.49\textwidth}@{\hfill}p{0.49\textwidth}@{}}
    Mengetahui KaProdi & Pembimbing Tugas Akhir \\[3cm]
    {\small\koorprodi} & {\small\bfseries\pembimbing} \\
\end{tabular}
}
